% ============================================================
% CHAPTER 9 — Discussion (FINAL restructure, further trimmed)
% ============================================================

\chapter{Discussion}
\label{chap:discussion}

This chapter discusses the results of Chapters~\ref{chap:results}--\ref{chap:external_validation}: why the agent behaves as it does, what the results mean for compliance-aware process optimisation, and which limitations remain, rather than repeating the experimental outcomes themselves. Four themes carry this argument. The agent recovers Pareto-optimal behaviour not because it is told which traces are optimal, but because a brevity-biased step penalty and the duration weight push it there together (Section~\ref{sec:discussion_tradeoff}). Compliance is never learned in the statistical sense: it is guaranteed by the shield from the first episode, and training instead reduces the agent's \emph{reliance} on that shield over time (Section~\ref{sec:discussion_compliance}). Role specialisation follows directly from which events have the largest cost--duration gap between resources, not from anything specific to reinforcement learning (Section~\ref{sec:discussion_roles}). And across three independent paradigms sharing no data, algorithm, or process formalism, the same cost-versus-duration trade-off keeps reappearing, the strongest evidence in this thesis that the trade-off is a property of the domain rather than of any one method (Section~\ref{sec:discussion_prior_work}). Section~\ref{sec:limitations} closes the chapter with its limitations; Chapter~\ref{chap:conclusion} then revisits the research questions and outlines future work.

% ============================================================
\section{How the Agent Learns the Cost--Duration Trade-off}
\label{sec:discussion_tradeoff}
\label{sec:discussion_reward}
\label{sec:discussion_reward_design}

The Expense Report experiment shows that the \ac{PPO} agent is able to discover Pareto-optimal process executions without explicitly searching through all possible traces. During training, the agent explored 524 different accepting traces within the $k{=}5$ limit, but only two of them were actually Pareto-optimal. Even so, the agent consistently converged towards both optimal solutions simply through interaction with the environment and the reward signal. This section first describes that convergence behaviour, then explains the reward-design mechanism that produces it.

\subsection{Recovering the Pareto front through reward weighting}

The two Pareto-optimal traces represent opposite optimisation goals: P1 is cheaper but slower, P2 is faster but more expensive. Which one a given weight configuration converges to (Section~\ref{sec:pareto_recovery}) follows directly from the brevity-bias mechanism described below.

\subsection{Suboptimal but valid traces}

A more interesting observation appears when analysing the full distribution of traces. Figure~\ref{fig:trace_distribution} shows that many episodes converge towards the trace:

\begin{center}
\texttt{FillOut $\to$ Approve $\to$ Withdraw}
\end{center}

This trace is fully valid according to the \ac{DCR} constraints, but it is still worse than both Pareto-optimal solutions in terms of cost and duration.

\begin{figure}[H]
    \centering
    \includegraphics[width=0.8\linewidth]{%
        Pictures/expense_trace_distribution.png}
    \caption{Distribution of accepting traces across all weight
             configurations during the final 20\% of training.
             The dominated trace
             \texttt{FillOut~$\to$~Approve~$\to$~Withdraw}
             appears frequently when duration is not penalised,
             but disappears once a duration penalty is introduced.}
    \label{fig:trace_distribution}
\end{figure}

This happens because after executing \texttt{Approve}, both \texttt{PayOut} and \texttt{Withdraw} are still legal actions. If duration is not part of the reward function, the agent has little reason to prefer one over the other. As a result, it can settle into traces that are compliant and locally acceptable, but globally suboptimal. Adding even a small duration penalty is enough to remove this behaviour almost entirely, since the shorter P2 trace becomes clearly more attractive.

The remaining category in Figure~\ref{fig:trace_distribution}, \emph{other valid traces}, stays a minor residual under $\beta=0$ (7--8\,\%) but rises under mixed weights, notably to 17.6\,\% at $\alpha=0.5,\,\beta=0.5$, consistent with the broader exploration sustained by balanced rewards (Section~\ref{sec:discussion_diversity}).


\paragraph{Key observation.} Compliance alone is not sufficient to guide the agent towards the best executions: the reward function must also distinguish between valid traces in a meaningful way. The next subsection identifies the specific mechanism, the step penalty, that does most of this work.

\subsection{The reward decomposition}

The reward function decomposes into additive components applied at each environment step (Table~\ref{tab:reward_components}, Equation~\ref{eq:reward_total}, Section~\ref{sec:reward}); combining them gives a net per-step effect of $-10$ for an illegal action, $+100$ on acceptance, $-0.5$ for a legal step with no progress, $+1.5$ for a step that resolves a pending event, and $-1.5$ for a repeated action, before the cost and duration penalties ($\alpha\hat{c}$ and $\beta\hat{d}$) are subtracted according to the current weight pair.

\paragraph{Brevity bias.} Because the step penalty accumulates unconditionally, shorter traces are always preferred over longer ones, all else being equal: a trace completed in two steps incurs a total step cost of $-3.0$, while a three-step trace incurs $-4.5$. The difference is modest in absolute terms but sufficient to influence policy gradients at the margin, and it is visible directly in the Expense Report results. Both Pareto-optimal traces are discovered under the duration-focused configuration ($\alpha=0, \beta=1$), not under the cost-focused one ($\alpha=1, \beta=0$): the cost-minimising configuration achieves a mean episode length of 3.4 steps and does not converge to Pareto-optimal traces, whereas the duration-minimising configuration achieves 2.3 steps on average and recovers both points, the short, expensive trace (cost $= 710$, duration $= 2{,}700$, two steps via \textit{WithdrawExpenseReport}) and the longer, cheap trace (cost $= 50$, duration $= 10{,}500$, three steps via \textit{Approve} $\to$ \textit{PayOut}). This occurs because the step penalty and the duration penalty act in the same direction under $\beta > 0$: both discourage long episodes, and the combined signal is strong enough that the agent also recovers the cheap three-step trace. The cost-focused agent, receiving no duration signal, relies solely on the step penalty to shape episode length and does not converge to the cost-optimal path within the training budget.

\paragraph{Step penalty versus terminal bonus.} The dominance of the terminal bonus means that the agent's primary objective is always to reach acceptance; the step penalty and the cost/duration components shape \emph{how} the agent reaches acceptance but do not override the imperative to reach it. This ordering is desirable: it ensures that the shield and the acceptance criterion remain the binding constraints, and that the cost/duration objectives function as secondary preferences rather than hard targets. A consequence, however, is that the reward signal provides relatively weak guidance for discriminating between traces of similar length. In the Expense Report, the two-step and three-step paths differ by only $1.5$ in accumulated step penalty, which leaves limited gradient headroom for the cost/duration components to differentiate the two Pareto points. The progress bonus ($+2$ per newly resolved pending event) and novelty penalty ($-1$ for repeated actions) partially counteract this by rewarding forward progress and discouraging cycles, but their effect is secondary to the terminal bonus in determining which trace the agent converges to.

\paragraph{Implications for processes with long mandatory traces.} For processes where the minimum compliant trace is inherently long, due to many mandatory events or complex response chains, the accumulated step penalty grows proportionally: a process requiring twenty mandatory steps incurs a step cost of $-30.0$ before any cost or duration penalty is applied. In such settings, the brevity bias may systematically suppress exploration of valid but lengthy execution paths, biasing the recovered Pareto front toward short, expensive traces and away from cost-optimal solutions that require more steps. Mitigating this is one of the directions identified in Section~\ref{sec:future_work}.

\subsection{Reward design as an implicit diversity mechanism}
\label{sec:discussion_diversity}

A related observation concerns the distribution of Pareto-optimal points across weight configurations in the Loan Application experiment. As reported in Section~\ref{sec:loanapp_pareto}, 15 of the 26 Pareto-optimal points originate from the balanced configuration ($\alpha=0.5,\,\beta=0.5$) rather than from the specialised agents. Under a pure scalarisation scheme one would expect specialised agents to dominate the interior of the front, with the balanced agent contributing only near the centre. Instead, the balanced reward appears to sustain broader trace exploration throughout training, preventing the premature convergence that specialised agents exhibit once their primary objective is satisfied.

This is consistent with the brevity bias described above: cost-focused agents converge quickly to short low-cost traces and stop exploring, while duration-focused agents converge to short fast traces. The balanced agent, receiving competing signals from both objectives simultaneously, cannot converge as cleanly to either extreme and continues exploring intermediate regions of the trace space. This suggests that, for this process, the balanced weight acts as an implicit diversity mechanism rather than merely interpolating between the two poles, though confirming whether this generalises to other \ac{DCR} graphs would require further experimentation.

This does not make the specialised configurations redundant. The balanced agent's exploration advantage is concentrated in the interior of the front: it is the specialised agents, not the balanced one, that consolidate the two extreme points, the cheapest and the fastest trace, which anchor the front and which a balanced reward has no incentive to converge to exclusively. The weight sweep is therefore not six attempts at the same goal with one winner; the balanced configuration covers breadth, and the specialised configurations cover the extremes that breadth alone does not reliably reach, which is why the full sweep, rather than a longer run of the balanced configuration alone, is needed to recover the front reported in Section~\ref{sec:loanapp_pareto}.

Despite the brevity bias, the normalised reward formulation proved sufficient to recover the exact Pareto front on the Expense Report benchmark and a well-populated approximation on the Loan Application, suggesting that the current design is adequate for processes of moderate complexity within the evaluated training budgets.

% ------------------------------------------------------------
\section{Compliance as a Learning Constraint}
\label{sec:discussion_compliance}

A common assumption in constrained optimisation is that limiting the agent's available actions reduces exploration and makes learning more difficult. The experiments suggest the opposite.

In the proposed framework, the \ac{DCR} shield blocks illegal actions before they can be executed, so the agent only interacts with process states that are actually reachable within the \ac{DCR} model; instead of corrupting the marking, an illegal attempt returns a fixed penalty that discourages repetition without affecting it. Chapter~\ref{chap:internal_validation}'s ablation (Table~\ref{tab:ablation_validity}) shows this empirically rather than conceptually: the unshielded agent spends its entire training budget exploring a corrupted state space and never recovers a single structurally valid accepting episode, while the shielded agent explores only valid states from the first episode. Rather than hurting optimisation, the shield helps stabilise it.

This is different from many safe \ac{RL} approaches in the literature, where illegal actions are discouraged only through large negative rewards~\cite{Guan2024Rule-basedControl, Odriozola-Olalde2025TowardsFramework, Gerold2026SafeShielding}. In those approaches, invalid actions may still occur during training. In the proposed framework, they are prevented directly by the \ac{DCR} semantics.

\subsection{The agent gradually stops relying on the shield}

Table~\ref{tab:compliance_milestones} shows how quickly the agent learns to avoid illegal actions on the Expense Report benchmark, the same phenomenon Chapter~\ref{chap:internal_validation} quantifies per training decile on the Loan Application benchmark.

\begin{table}[H]
\centering
\caption{Illegal-action milestones during training
         (Expense Report benchmark, normalised reward).}
\label{tab:compliance_milestones}
{\renewcommand{\arraystretch}{1.5}
\setlength{\tabcolsep}{14pt}
\begin{tabular}{lc}
\toprule
\textbf{Milestone} & \textbf{Value} \\
\midrule
Illegal ratio below 50\% & 580 steps ($0.6\,\%$ of training budget) \\
Final illegal ratio (across configurations) & $0.2\,\%$--$10.0\,\%$ \\
Acceptance rate (all configurations) & $100\,\%$ \\
\bottomrule
\end{tabular}}
\end{table}

The illegal-action ratio drops below 50\% after only 580 steps, which corresponds to less than 1\% of the total training budget. By the end of training, four of the six configurations fall below 3\%, while the remaining two, the pure cost and pure duration minimisers, stabilise around 10\%. This suggests that the shield is not simply forcing compliance from outside the learning process. Instead, the agent gradually learns the structure of the \ac{DCR} graph and begins to avoid invalid actions on its own.

A similar pattern appears in the Loan Application experiment, even though the action space expands from 12 to 24 actions (one per event--role pair) relative to the flat LoanApp variant. The agent still converges towards highly compliant behaviour, suggesting that the shielding mechanism scales well with increasing action-space complexity, at least up to the 24-action range tested here; Chapter~\ref{chap:external_validation} extends this question to much larger action spaces (up to 113 events) and finds that this scaling does not continue indefinitely without additional support such as action masking (Section~\ref{sec:external_discussion}).

% ------------------------------------------------------------
\section{Role Specialisation and Event-level Granularity}
\label{sec:discussion_roles}

The patterns observed in Section~\ref{sec:loanapp_roles} reflect a clear dependency between optimisation objectives and role selection. This section interprets the structural reasons behind this behaviour; Section~\ref{sec:discussion_prior_work} situates it within prior work and identifies what the proposed framework contributes that existing approaches cannot provide.

\subsection{Why sensitivity differs across events}

The variation in slope magnitude in Figure~\ref{fig:loanapp_role_slopes} can be explained by how the reward function scales cost and duration. Since both components are normalised by their maximum values in the process graph, the effective impact of switching between Junior and Expert depends on the base cost and duration of each event. High-cost activities such as \textit{AML check} (base cost 539) and \textit{Appraise property} (base cost 509) produce large absolute reward differentials between roles, generating a strong policy gradient signal and driving near-complete specialisation at both weight poles. Lower-impact activities produce smaller differentials, and the gradient signal is insufficient to consistently override exploration noise.

\textit{Assess loan risk} represents a boundary case: its cost--duration profile places it close to the indifference threshold of the normalised reward, resulting in approximately 50\,\% Expert assignment regardless of weight configuration. This is actionable from a process management perspective: role-indifferent activities can be staffed on availability grounds without measurable impact on the cost--duration trade-off, while optimisation effort should concentrate on the high-sensitivity events identified by steep slopes in Figure~\ref{fig:loanapp_role_slopes}.

\paragraph{Key observation.} Role specialisation is not an emergent RL phenomenon so much as a direct readout of which events have the largest cost--duration gap between resources. The framework's contribution is exposing this gap at event-level granularity, which Section~\ref{sec:discussion_prior_work} contrasts with the resource-pool granularity of prior approaches, and which Chapter~\ref{chap:conclusion} returns to as contribution~C3 of the thesis.

% ------------------------------------------------------------
\section{Comparison with Prior Work}
\label{sec:discussion_prior_work}

This section places the results alongside exact constraint solving and two independent lines of prior multi-objective process optimisation work. Cost and duration values are not directly comparable in absolute terms across approaches: Kirchdorfer et al.~\cite{KirchdorferLukasFromPolicies} report labour costs in dollars and waiting time in hours derived from real event logs, Doumeni~\cite{ArtemisDoumeniParetoAgentSimulator} reports man-hour cost and waiting time relative to an As-Is simulation baseline, while this work uses abstract units with fixed role multipliers. Comparisons should therefore be interpreted in terms of directional and structural similarity rather than numerical equivalence.

\subsection{Exact Constraint Solving}
\label{sec:discussion_csp_comparison}

\paragraph{Key observation.} Both the \ac{PPO} agent and the CSP+COP solver recover exactly the same Pareto front on the Expense Report benchmark. The difference lies mainly in how the solutions are found and what guarantees each method provides.

Table~\ref{tab:method_comparison} summarises the main differences between the two approaches.

\begin{table}[H]
\centering
\caption{Conceptual comparison between CSP+COP and PPO optimisation
         on the Expense Report benchmark.}
\label{tab:method_comparison}
{\renewcommand{\arraystretch}{1.5}
\begin{tabular}{p{3.5cm} p{5cm} p{5cm}}
\toprule
\textbf{Aspect} & \textbf{CSP+COP} & \textbf{PPO (proposed)} \\
\midrule
Optimisation strategy &
Exhaustive symbolic search &
Learning through interaction \\

\addlinespace
Pareto guarantee &
Exact &
Approximate \\

\addlinespace
Compliance handling &
Constraint satisfaction &
Shielded execution \\

\addlinespace
Runtime on benchmark &
1.01\,s &
$\sim$82\,min training (full sweep) \\

\addlinespace
Main scalability issue &
Constraint explosion &
Training budget and exploration \\

\addlinespace
Required input &
Symbolic optimisation model &
DCR graph + reward function \\

\bottomrule
\end{tabular}}
\end{table}

The CSP+COP method guarantees that no Pareto-optimal solution is missed within the chosen search bound, a property the \ac{RL} approach cannot replicate by design. What the results demonstrate, however, is that this gap matters less than it might appear: for processes with a small and convex Pareto front, the sampling-based approximation converges to the exact solution within a moderate training budget. The more important question for practical deployment is therefore not whether the front is exact, but whether the compliance guarantee and the training cost are acceptable trade-offs for the scalability gained.

\paragraph{Training cost versus inference cost.} The 1.01\,s versus $\sim$82\,min comparison in Table~\ref{tab:method_comparison} is misleading if read as a per-query cost, because the two numbers measure different things. The CSP+COP solver re-solves an NP-hard constraint optimisation problem from scratch every time a trace is requested, so 1.01\,s is the marginal cost of every single query. The $\sim$82\,min figure for \ac{PPO} is the one-time cost of training all six weight configurations (mean $13.6$\,min per configuration, measured from consecutive run timestamps), amortised over the lifetime of the resulting policies: once trained, generating a new trace is a single forward pass through the policy network, with no search involved. Measured directly on the trained Expense Report policy (duration-focused configuration, $\alpha{=}0,\beta{=}1$), a single inference call takes $0.08$\,ms on average (2{,}000 calls, CPU), a $\sim$12{,}400$\times$ reduction relative to the solver's 1.01\,s per query, and is effectively constant regardless of how many traces are requested afterwards. The comparison that matters for deployment is therefore not solver time versus training time, but solver time versus \emph{inference} time once training is paid for once; viewed this way, the \ac{RL} approach is the cheaper option for any deployment that needs to generate traces repeatedly, such as recommending an action to a process participant in real time, while the exact solver remains preferable when a single optimal answer is needed once and the process is small enough for it to converge at all.

\paragraph{What the forward pass actually computes.} The 0.08\,ms figure is not a black box: it is the cost of one evaluation of the 64-unit, two-hidden-layer policy network of Section~\ref{sec:ppo_algorithm} on the 15-dimensional Expense Report observation of Section~\ref{sec:state}, no search, no constraint check, and no call to the \ac{DCR} engine, which is the structural reason it is several orders of magnitude faster than a constraint solver invocation. It also explains why the policy network alone cannot be the compliance mechanism: at the initial marking, the highest-scoring event under the trained weights is \textit{Approve}, which is not enabled, since \textit{FillOut} has not been executed yet. The network was never required to learn this constraint, because the shield (Section~\ref{sec:discussion_compliance}) intercepts the proposal regardless of what the policy outputs; compliance is enforced by the architecture, not by the accuracy of the forward pass.

This is the worked example promised in Section~\ref{sec:training_pipeline} of the policy/trace distinction: the policy here is the fixed vector of all $5{,}509$ weights and biases of the actor network just evaluated; the trace is the sequence of executed events produced by repeatedly applying it within one episode. The same trained policy yields different traces depending on where an episode starts and which proposals the shield happens to intercept along the way, which is why the Pareto fronts reported in Chapter~\ref{chap:results} are computed over the traces a policy produces, not over its parameters directly.

\paragraph{Scalability trade-off.} The two approaches scale differently as process complexity increases. In the symbolic approach, performance is mainly affected by the number of constraints and the complexity of the optimisation problem; as the process grows, the search space becomes increasingly difficult to explore exhaustively. This is no longer a purely theoretical concern: Chapter~\ref{chap:external_validation} measures exactly where this happens. On the mined-graph grid, the CSP+COP solver stops converging within a 300\,s budget somewhere between 42 and 113 activities, and the \ac{PPO}-based agent's accept rate collapses within that same wide interval (Section~\ref{sec:external_scalability}); a finer-grained diagnostic on a separate, controlled graph series narrows this shared crossover to between 40 and 60 activities (Section~\ref{sec:external_scaling_diagnostic}), for a structurally different reason in each method: not constraint explosion for \ac{PPO}, but the size of the action space the policy must search without guidance.

The \ac{RL} framework has a different limitation: training becomes more expensive when the action space grows or when more interaction steps are required. For example, introducing role-conditioned actions in the Loan Application process doubles the number of available actions from 12 to 24; even so, the agent is still able to converge towards a well-populated Pareto front, well within the range where the learning component still adds value over a purely structural baseline (Section~\ref{sec:discussion_roles}), though this should not be read as evidence that the framework scales past where exact enumeration becomes infeasible. Another advantage is that compliance handling is separated from the learning objective itself: the shielding mechanism blocks invalid actions directly at the environment level, meaning the agent does not need to learn process constraints purely through punishment.

Overall, the two approaches are complementary rather than competing, but not because \ac{RL} scales further in terms of process size: at no point in this thesis, on either graph series, does \ac{RL} recover a valid front on a graph where the exact solver could not. The complementarity is instead in what each method guarantees once it does work: CSP+COP guarantees an exact, exhaustive Pareto front within its search budget, while \ac{RL} trades that exactness for an approximate front obtained through interaction rather than symbolic search, an advantage in setup cost and in the wall-clock comparison with population-based evolutionary search of Section~\ref{sec:discussion_doumeni}, but not in the maximum process size either method can handle. The only condition in this study that scales past the point where exact enumeration becomes infeasible is the shielding mechanism on its own, without a learning component (Sections~\ref{sec:external_scalability}, \ref{sec:external_scaling_diagnostic}); within the size range both methods can solve, the value \ac{RL} adds is the quality of the trade-off it recovers, not the size of the process it reaches.

\subsection{Cross-Paradigm Convergence}
\label{sec:discussion_cross_paradigm}

Table~\ref{tab:cross_paradigm} places the results alongside three independent lines of prior work, summarising the cost--duration trade-off structure and compliance guarantees across all four approaches. The figures for this work use the capacity-constrained Loan Application results at budget $k=3$ (Section~\ref{sec:loanapp_capacity}), rather than the unlimited-capacity variant: under unlimited Expert availability role selection is separable across events and does not constitute a genuine allocation decision, whereas Kirchdorfer et al.\ and Doumeni both optimise resource allocation under realistic scarcity, making the constrained variant the closer comparison in spirit. A more detailed comparison between the present framework and Doumeni is provided in Section~\ref{sec:discussion_doumeni}.

\begin{table}[H]
\centering
\caption{Qualitative comparison of cost--duration trade-off
structure and compliance guarantees across three independent
approaches on the Loan Application domain. Figures for this work are
under the capacity-constrained variant ($k=3$).}
\label{tab:cross_paradigm}

{\renewcommand{\arraystretch}{1.5}
\begin{tabularx}{\linewidth}{X c c p{3.2cm}}
\toprule
\textbf{Approach} & \textbf{Cost} &
\textbf{Duration} & \textbf{Compliance} \\
\midrule
Kirchdorfer et al.\ (NSGA-II)
    & $\sim -32\%$ & strong reduction
    & Statistical \\
Doumeni (NSGA-II)\textsuperscript{b}
    & $\sim -16\%$ & $\sim -10\%$
    & Statistical \\
Huang et al.\textsuperscript{a}
    & task-level & task-level
    & Not addressed \\
This work ($\alpha=1, \beta=0$, $k=3$)
    & $\sim -22\%$ & $+22\%$
    & By construction \\
This work ($\alpha=0, \beta=1$, $k=3$)
    & $+41\%$ & $\sim -30\%$
    & By construction \\
\bottomrule
\multicolumn{4}{l}{%
  \footnotesize\textsuperscript{a}~Relative improvements over
  baselines only; no absolute values published.}\\
\multicolumn{4}{l}{%
  \footnotesize\textsuperscript{b}~Reported vs.\ As-Is baseline; not directly comparable to RL baseline.}
\end{tabularx}}

\vspace{2mm}

\end{table}

The three approaches share no data, no algorithm, and no process formalism, yet all converge to the same qualitative conclusion: minimising execution cost requires slower, cheaper resources, and minimising duration requires faster, more expensive ones, with no single assignment dominating both objectives simultaneously. This cross-paradigm convergence suggests that the trade-off structure is a genuine economic property of resource allocation in business processes rather than an artefact of any particular optimisation method.

Beyond this shared structure, the present analysis contributes three properties that distinguish it from the other approaches in this comparison; Chapter~\ref{chap:conclusion} states these formally as contributions C1--C3 of the thesis.

\subsection{Comparison with Evolutionary Multi-objective Search}
\label{sec:discussion_doumeni}

Doumeni~\cite{ArtemisDoumeniParetoAgentSimulator} addresses the same optimisation problem studied in this thesis: balancing operational cost and execution time in the Loan Application domain. However, the two approaches rely on fundamentally different modelling and optimisation paradigms. Doumeni applies a multi-objective Genetic Algorithm (NSGA-II) to agent assignment policies within the AgentSimulator framework, whereas this work formulates the problem as reinforcement learning over a declarative \ac{DCR} process model. Table~\ref{tab:doumeni_comparison} expands on the Doumeni entry in Table~\ref{tab:cross_paradigm}, detailing the comparison across six dimensions before discussing each aspect in turn.

\begin{table}[H]
\centering
\caption{Comparison between the NSGA-II approach of Doumeni and the proposed DCR-based reinforcement learning framework.}
\label{tab:doumeni_comparison}

{\renewcommand{\arraystretch}{1.5}

\begin{tabularx}{\linewidth}{l p{4.2cm} p{4.2cm}}
\toprule
\textbf{Dimension} & \textbf{Doumeni (NSGA-II)} & \textbf{This Work} \\
\midrule

Optimisation method &
Multi-objective evolutionary search &
Reinforcement learning (PPO) \\

Process representation &
AgentSimulator process model learned from event logs &
Declarative DCR process model \\

Compliance mechanism &
Inherited from simulation behaviour and historical data &
Formally enforced through DCR shielding \\

Input requirements &
Historical event logs &
Normative process model \\

Decision granularity &
Resource pools and process stages &
Event-level activity and role assignment \\

Multi-objective handling &
Pareto search through population evolution &
Weight-conditioned reward $(\alpha,\beta)$ \\

\bottomrule
\end{tabularx}}
\end{table}

Figure~\ref{fig:approach_comparison} illustrates the structural difference between the two pipelines. Both follow a four-stage architecture, input, process representation, optimisation method, and output, but differ fundamentally in their starting point and compliance mechanism, as discussed below.

\begin{figure}[H]
    \centering
    \includegraphics[width=\linewidth]{Pictures/thesis_approach_artemis_comparison.png}
    \caption{Structural comparison between the NSGA-II approach
of Doumeni~\cite{ArtemisDoumeniParetoAgentSimulator} (top)
and the proposed DCR-based reinforcement learning framework
(bottom). Both pipelines produce a Pareto set of policies
optimising cost and duration, but differ in process
representation, optimisation method, and the nature of the
compliance guarantee.}
    \label{fig:approach_comparison}
\end{figure}

\paragraph{Structural convergence.} The cost--duration trade-off recovered here follows the same direction as Doumeni's, already quantified side by side in Table~\ref{tab:cross_paradigm}: improvements in one objective come at the expense of the other, and no single policy dominates both simultaneously, despite the two methods sharing no data, algorithm, or process formalism.

\paragraph{Scalability.} A practical distinction concerns computational cost. The structural reason is that evolutionary search evaluates a population of policies across multiple generations, each requiring hundreds of full process simulations, whereas the \ac{RL}-based approach amortises this cost across training steps rather than across population members. Each $(\alpha,\beta)$ configuration is directly comparable to one NSGA-II run in terms of what it produces, a set of policies covering a region of the trade-off space. On a per-configuration basis, training time ranges from $9.4$ to $11.8$\,min across the six weight configurations (mean $10.1$\,min, measured from consecutive run timestamps),\footnote{This is slightly lower than the $13.6$\,min per configuration measured on the smaller Expense Report benchmark (Section~\ref{sec:discussion_csp_comparison}), despite Loan Application's larger graph. The reason is that wall-clock training time has two additive components: a per-step cost shared by both benchmarks (the same $100{,}352$ steps, dominated by the environment's network round-trip rather than graph size) and a per-episode reset cost. Because Expense Report accepts in only $2$--$3$ steps, it completes roughly three times more episodes within the same step budget ($42{,}351$ versus $14{,}563$), and therefore pays the reset cost more often; fitting both measurements to this two-component model gives a per-step cost of $4.9$\,ms and a per-reset cost of $7.6$\,ms, both plausible for a local HTTP round-trip to the Node.js adapter.} a roughly $33\times$ reduction in wall-clock time against Doumeni's $339$\,min per NSGA-II run, and the complete six-configuration weight sweep, which approximates the full Pareto front, completes in about one hour. This architectural difference suggests that the \ac{RL} framework may be better suited to settings where simulation-per-evaluation becomes prohibitive, though whether this advantage holds for significantly larger \ac{DCR} graphs remains to be verified experimentally. It is worth noting that this advantage is specific to the comparison with population-based evolutionary search: it does not extend to the comparison with the exact CSP+COP solver of Section~\ref{sec:discussion_csp_comparison}, which is faster still on the small graphs where it converges, nor does it close the action-space ceiling identified in Chapter~\ref{chap:external_validation}.

\paragraph{Compliance guarantees.} Another important distinction concerns process validity. AgentSimulator generates traces by sampling from probabilistic transition models derived from historical event logs; compliance is therefore tied to the behaviour represented in the data and preserved through simulation fidelity, but the framework does not provide a mechanism for verifying whether newly generated traces satisfy an explicit process specification. The present framework operates under a different assumption: actions are filtered through the \ac{DCR} execution semantics before they can be executed, ensuring that every generated trace satisfies the declarative process constraints. This property holds throughout training and execution rather than emerging only after optimisation has converged.

\paragraph{Process representation.} The two approaches should therefore be viewed as complementary rather than competing alternatives. Doumeni follows a data-driven perspective, learning process behaviour, resource schedules, and transition probabilities from historical event logs before optimisation takes place, which makes the approach particularly suitable when reliable operational data is available but no formal process specification exists. The present framework starts from the opposite direction: a declarative \ac{DCR} model acts as the normative description of allowed behaviour, and reinforcement learning is used to discover resource-allocation policies within that space, with the resulting compliance guarantees stemming directly from this modelling choice.

% ------------------------------------------------------------
\section{Limitations}
\label{sec:limitations}

The following limitations constrain the scope of the conclusions that can be drawn from the present work, grouped by whether they concern a methodological choice that was not explored further, a measurement that was not carried out, or a gap between the model and a real deployment. Section~\ref{sec:future_work} discusses how several of them could be addressed.

\paragraph{Methodological choices not explored further.}

\begin{itemize}

    \item Single RL algorithm.
    All experiments use PPO as the sole learning algorithm.
    Whether other policy gradient or value-based methods would
    produce different Pareto coverage or convergence behaviour
    on the same benchmarks has not been evaluated.

    \item No systematic hyperparameter search.
    Hyperparameters other than the scalarisation weights
    $(\alpha,\beta)$ themselves (Table~\ref{tab:hyperparams}) were
    fixed from preliminary runs and shared across both benchmarks,
    rather than tuned per benchmark or searched systematically. This
    is a threat to internal validity: the reported results establish
    that the chosen configuration works, not that it is optimal, and
    a different configuration could plausibly change convergence
    speed or Pareto coverage.

    \item Linear scalarisation.
    The multi-objective optimisation relies on linear
    scalarisation of cost and duration.
    This approach cannot recover Pareto points in non-convex
    regions of the objective space.
    For the benchmarks used here the front is convex and
    the limitation does not affect the reported results,
    but it may become relevant for processes with more
    complex trade-off geometries.

    \item Truncation handled as ordinary termination.
    Episodes that exhaust the step limit (\texttt{MAX\_EPISODE\_STEPS})
    without reaching an accepting marking are reported to the learner
    as an ordinary termination rather than a Gymnasium-style
    truncation, so the value function never bootstraps from the
    cutoff state as it ideally should. Given the $100\,\%$ acceptance
    rates reported throughout most of this thesis, such episodes are
    rare in practice, but the simplification was not corrected, nor
    was its impact on the value estimates quantified.

\end{itemize}

\paragraph{Measurements not carried out.}
\begin{itemize}

    \item Seed robustness was measured on Loan Application, not on Expense Report.
    The headline weight sweep was repeated under $n=4$ random seeds on
    the Loan Application benchmark, with stable results
    (Section~\ref{sec:policy_generalisation}); the Expense Report
    benchmark was run only under the reported seed ($s=1$). Transfer
    across structurally different process graphs, rather than across
    seeds of the same graph, is addressed separately in
    Chapter~\ref{chap:external_validation}.

    \item Approximated Pareto front.
    The Pareto front is extracted from sampled accepting
    episodes rather than computed exactly.
    Low-probability Pareto-optimal traces may not appear
    in the sample and would therefore be missed.
    For the Expense Report benchmark, the two-point ground
    truth is fully recovered; for the Loan Application
    variant with roles, no exact reference exists and the
    reported front may be incomplete.

\end{itemize}

\paragraph{Gap between the model and a real deployment.}
\begin{itemize}

    \item Brevity bias from the step penalty.
    The unconditional per-step penalty that shapes episode length
    systematically favours shorter compliant traces over longer ones,
    all else being equal; for processes whose minimum compliant trace
    is inherently long, this may suppress exploration of valid but
    lengthy execution paths and bias the recovered front away from
    cost-optimal solutions that require more steps, as discussed in
    Section~\ref{sec:discussion_tradeoff}. This did not affect the
    benchmarks evaluated here.

    \item Synthetic cost models.
    Cost and duration values are taken from benchmark
    definitions or assigned via fixed role multipliers.
    Real processes exhibit stochastic and context-dependent
    resource costs not modelled here, which limits the
    direct applicability of the quantitative results to
    industrial settings.

    \item No continuous time or concurrency.
    Duration is modelled as a per-event deterministic scalar.
    Resource contention, waiting times, and concurrent
    activity execution are not represented, in contrast to
    simulation-based approaches such as
    Doumeni~\cite{ArtemisDoumeniParetoAgentSimulator}.

    \item Action-space scalability without masking.
    As discussed in Section~\ref{sec:discussion_csp_comparison}, the
    unmasked \ac{PPO} agent's accept rate collapses beyond
    approximately 42 activities because the policy still
    samples over the full action set rather than only the currently
    enabled subset; only the shielding mechanism alone, without
    learning, continues to scale past that point.

    \item Acceptance rate as compliance proxy.
    Acceptance in the \ac{DCR} semantics is a well-defined
    formal property, but it does not capture all possible
    notions of process correctness such as data constraints,
    real-time deadlines, or resource availability.

\end{itemize}

Chapter~\ref{chap:conclusion} synthesises the themes of this chapter into direct answers to the thesis's research questions, restates its contributions, and turns the limitations above into concrete directions for future work.
